% Generated draft: requires mathematical and source review.
\documentclass[11pt]{article}
\usepackage[margin=25mm]{geometry}
\usepackage{fontspec}
\setmainfont{Latin Modern Roman}
\usepackage{amsmath,amssymb,mathtools,unicode-math}
\setmathfont{Latin Modern Math}
\usepackage{enumitem,needspace,xurl,hyperref}
\setlength{\parindent}{0pt}
\setlength{\parskip}{5pt}
\setlength{\emergencystretch}{3em}
\newcommand{\N}{\mathbb{N}}
\newcommand{\Z}{\mathbb{Z}}
\newcommand{\Q}{\mathbb{Q}}
\newcommand{\R}{\mathbb{R}}
\newcommand{\C}{\mathbb{C}}
\newcommand{\F}{\mathbb{F}}
\newcommand{\A}{\mathbb{A}}
\newcommand{\PP}{\mathbb{P}}
\newcommand{\E}{\mathbb{E}}
\newcommand{\eps}{\varepsilon}
\newcommand{\dd}{\,\mathrm{d}}
\begin{document}
\section*{UnsolvedMath ID: 30001690}
\textbf{Source problem code:} catalogue entry
\subsection{Definitions and mathematical statement}
Let \(V\) and \(E\) denote the vertex and edge sets of the discrete geometric object intended by the source. An edge-length assignment is a function \(\ell:E\to\R_{>0}\), and an assignment of positive vertex scale factors is a function \(s:V\to\R_{>0}\). These notations specify indexing and positivity only; the source does not specify the underlying discrete geometric objects, their dimension, their topology, or their admissibility conditions.

The source describes \emph{discrete conformal equivalence} as a relation in which corresponding edge lengths differ by positive scale factors attached to vertices. It does not state the formula relating an edge length to the scale factors at its endpoints. Consequently, no particular vertex-scaling formula is imposed here. The source also does not define a \emph{spin transformation}, its domain, or the conditions under which it produces an admissible transformed discrete object.

\textbf{Problem as supplied.} Can this notion of discrete conformal equivalence be extended to all spin transformations? The requested extension should recover the stated vertex-scaling notion wherever that notion applies and encompass every spin transformation in the intended class.

\textbf{Unresolved formulation qualifications.} The intended class of discrete objects, the exact vertex-scaling relation, the definition and admissibility conditions of spin transformations, and the mathematical meaning of the proposed extension must be specified before this question determines a fully precise mathematical assertion. In particular, the quantifier ``all spin transformations'' is retained, but its domain cannot be recovered from the supplied record. No additional hypotheses, boundary conditions, or extension axioms are assumed.
\subsection{Short English statement}
The problem asks whether discrete conformal equivalence, described through positive scale factors at vertices, can be extended to every spin transformation. The supplied record does not define the scaling rule, the discrete objects, or spin transformations sufficiently to make the requested extension precise. This is a statement-only draft requiring human mathematical review.
\subsection{Sources}
\begin{itemize}
\item Local UnsolvedMath registry: supplied complete record for permanent ID \(30001690\), titled \emph{Spin Transformations of Discrete Conformal Structures}, source problem code \texttt{catalogue entry}, in the Oberwolfach Reports Open Problems collection. Project registry locations are \texttt{data/unsolvedmath/problems.json} and \texttt{catalog/problems.json}; these files were not independently inspected for this draft.
\item Supplied source citation: \emph{Trends in Mathematical Imaging and Surface Processing} (2011), Oberwolfach Reports, pp.\ 293--348, DOI \texttt{10.4171/owr/2011/07}. Original report: \url{https://doi.org/10.4171/owr/2011/07}. The record attributes the question to workshop open-problem discussions and supplies no individual proposer.
\item Additional supplied literature source, labeled ``discrete extrinsic and intrinsic Dirac operators'': \url{https://doi.org/10.1080/10586458.2020.1727798}. No authors, complete bibliographic citation, or specific supporting result are supplied for this reference.
\item Literature assessment dated exactly \texttt{2026-08-21}: the supplied curation labels the problem \emph{partially solved} and states that discrete Dirac and spin-transformation theory has advanced, while a general extension of the stated discrete conformal-equivalence notion to all spin transformations remains open. This is supplied dated curation, not an independently verified literature assessment; the record expressly requests independent verification of time-sensitive status and cited claims.
\end{itemize}
\end{document}
